% =====================================================================
%  1bat-ud4-14.tex  —  SESSIÓ 14: Repàs global i síntesi (estacions)
%  (sense preàmbul: s'inclou des de main.tex)
% =====================================================================
\horatitol{Sessió 14 — Repàs global i síntesi (estacions)}%
{Repàs actiu de tota la unitat per quatre estacions, una per bloc, i un mapa final que connecta cada eina amb la realitat social que descriu.}

\subsection{Idea clau}

Avui repassem \textbf{tota la unitat} de manera activa. Hi ha \textbf{quatre
estacions}, una per cada bloc que hem treballat. Rotareu per totes amb temps
controlat (uns $10$ minuts cada una), de manera que tornareu a tocar tots els
continguts en contextos una mica diferents. Al final, construirem entre tots un
\textbf{mapa de la unitat}.

\begin{idea}
\textbf{Les quatre estacions}
\begin{enumerate}[leftmargin=1.6em, itemsep=2pt]
  \item \textbf{Paràmetres:} calcular mitjana, mediana i desviació típica d'un conjunt
        i interpretar-ne la dispersió.
  \item \textbf{Lectura de gràfics:} llegir una piràmide de població o un histograma i
        treure'n una conclusió.
  \item \textbf{Lectura crítica:} detectar el parany d'un titular o un gràfic
        enganyós.
  \item \textbf{Enquestes:} revisar o redactar una pregunta neutra i tabular un petit
        conjunt de respostes.
\end{enumerate}
\textbf{Consell:} a cada estació, abans de calcular, pregunta't \emph{què} t'està
demanant el bloc i \emph{quina} realitat social hi ha al darrere.
\end{idea}

\subsection{Les estacions}

\begin{exemple}[Estació 1 — paràmetres]
Un conjunt d'edats d'usuaris: $20$, $22$, $21$, $19$, $48$. Calcula la mitjana i la
mediana, i digues quina representa millor el grup i per què.
\end{exemple}

\begin{exemple}[Estació 2 — lectura d'una piràmide]
Mira aquesta piràmide de població. Identifica la franja més nombrosa (moda) i digues
si el territori és jove o envellit.

\begin{center}
\begin{tikzpicture}
\begin{axis}[
    width=10cm, height=5cm,
    xbar, bar width=7pt, y=12pt,
    xmin=-18, xmax=18,
    xtick={-16,-8,0,8,16}, xticklabels={16,8,0,8,16},
    xticklabel style={font=\scriptsize, color=udgrey},
    xlabel={\small \% població \quad (homes $\leftarrow$ \ $\rightarrow$ dones)},
    ytick={1,2,3,4,5,6},
    yticklabels={0-14,15-29,30-44,45-64,65-79,80+},
    yticklabel style={font=\scriptsize, color=udgrey},
    tick style={udgrey}, axis line style={udgrey},
    enlarge y limits=0.12,
    legend style={font=\scriptsize, at={(0.5,-0.32)}, anchor=north, legend columns=2}]
  \addplot[fill=udblue!55, draw=udnavy] coordinates {(-14,1) (-15,2) (-13,3) (-9,4) (-5,5) (-2,6)};
  \addplot[fill=udblue!25, draw=udnavy] coordinates {(13,1) (14,2) (13,3) (9,4) (6,5) (3,6)};
  \legend{homes, dones}
\end{axis}
\end{tikzpicture}

\figcap{Estació 2: una piràmide per llegir. Quina és la moda? Jove o envellida?}
\end{center}
\end{exemple}

\subsection{Exercicis resolts}

\begin{exresolt}[model de l'Estació 1: paràmetres]
Per a les edats $20$, $22$, $21$, $19$, $48$, calcula la mitjana i la mediana i digues
quina representa millor el grup.
\solucio
\textbf{Mitjana:}
\[
\overline{x}=\frac{20+22+21+19+48}{5}=\frac{130}{5}=26.
\]
\textbf{Mediana:} ordenem $19,20,\underline{21},22,48$; la central ($N=5$) és
$\text{Me}=21$.

El valor $48$ (segurament un educador) \textbf{estira la mitjana} fins a $26$, però la
mediana ($21$) reflecteix molt millor un grup que és, de fet, de joves d'uns $20$
anys.
\resposta{$\overline{x}=26$, $\text{Me}=21$; la mediana representa millor el grup (el
$48$ distorsiona la mitjana).}
\end{exresolt}

\begin{exresolt}[model de l'Estació 3: lectura crítica]
Un titular diu: «la participació al casal ha crescut un $100\%$!». Quina pregunta
clau faries abans de celebrar-ho?
\solucio
La pregunta clau és «\textbf{$100\%$ de què?}», és a dir, quina era la \textbf{base}.
Un $100\%$ d'augment vol dir que s'ha \emph{doblat}: si abans hi havia $4$ infants,
ara n'hi ha $8$ (un canvi petit en termes absoluts); si abans n'hi havia $200$, ara
$400$ (un canvi enorme). Sense la base, el percentatge sol no diu la magnitud real.
\resposta{Cal preguntar la base: «$100\%$ de què?». Doblar de $4$ a $8$ no és el
mateix que de $200$ a $400$.}
\end{exresolt}

\subsection{Exercicis per fer}

\noindent\textit{Una tasca per estació. Rota i resol-les totes.}

\begin{exercicis}
\begin{exlist}
  \item \textbf{(Estació 1)} Per a les hores de voluntariat $3$, $5$, $4$, $5$, $3$:
        calcula la mitjana, la moda, la mediana i la desviació típica.
  \item \textbf{(Estació 2)} De la piràmide de l'exemple de l'Estació 2: quina és la
        franja moda? És un territori jove o envellit? Quins serveis hi prioritzaries?
  \item \textbf{(Estació 3)} «El barri A té $40$ casos i el B, $20$: el A és el doble
        de problemàtic.» Si A té $\num{8000}$ habitants i B en té $\num{2000}$, calcula
        la taxa per mil de cadascun. Qui en té més, proporcionalment?
  \item \textbf{(Estació 4)} Revisa aquesta pregunta i, si cal, fes-la neutra: «No
        trobes que el casal hauria d'obrir més hores?». Després, tabula aquestes $10$
        respostes a una pregunta de Sí/No: Sí, Sí, No, Sí, No, Sí, Sí, No, Sí, Sí.
  \item \textbf{(Mapa de la unitat)} Completa, amb les teves paraules, la connexió de
        cada eina amb la realitat social que descriu:
        \begin{center}
        \renewcommand{\arraystretch}{1.4}
        \begin{tabular}{p{4.8cm} p{7.8cm}}
        \toprule
        \textbf{Eina} & \textbf{Quina realitat social descriu / per a què serveix} \\
        \midrule
        Mitjana, mediana, moda & \rule{7.2cm}{0.4pt} \\
        Desviació típica & \rule{7.2cm}{0.4pt} \\
        Piràmide / histograma & \rule{7.2cm}{0.4pt} \\
        Lectura crítica & \rule{7.2cm}{0.4pt} \\
        Enquesta i buidatge & \rule{7.2cm}{0.4pt} \\
        \bottomrule
        \end{tabular}
        \end{center}
  \item \textbf{(Síntesi)} En una frase per bloc, resumeix què has après a la unitat:
        descriure dades, visualitzar-les, llegir-les críticament i recollir-ne de
        pròpies.
\end{exlist}
\end{exercicis}

% ---------------------------------------------------------------------
\fullsolucions{Sessió 14}
\begin{solucions}
\begin{sollist}
  \item Ordenat: $3,3,4,5,5$. Mitjana $\overline{x}=\dfrac{20}{5}=4$; moda bimodal
        ($3$ i $5$); mediana (central) $\text{Me}=4$. Distàncies al quadrat respecte de
        $4$: $1,1,0,1,1$, suma $4$; $\sigma^2=\dfrac{4}{5}=0,8$;
        $\sigma=\sqrt{0,8}\approx 0,89$.
  \item La moda és la franja \textbf{$15$–$29$ anys} i la base ($0$–$14$) és ampla: és
        un territori \textbf{jove}. Hi prioritzaria llars d'infants i escoles, casals i
        ocupació juvenil.
  \item Taxa A: $\dfrac{40}{\num{8000}}\per\num{1000}=5$ per mil. Taxa B:
        $\dfrac{20}{\num{2000}}\per\num{1000}=10$ per mil. \textbf{El barri B} en té
        més, proporcionalment ($10$ contra $5$ per mil), al contrari del que deia el
        titular.
  \item La pregunta original és \textbf{esbiaixada} («No trobes que\dots»); neutra:
        «Quantes hores creus que hauria d'obrir el casal?» o «Com valores l'horari
        actual del casal?» (Insuficient / Adequat / Excessiu). Tabulació de les $10$
        respostes: \textbf{Sí $=7$} ($70\%$), \textbf{No $=3$} ($30\%$).
  \item Resposta oberta. Idees esperades: \emph{centralització} $\rightarrow$ resumir
        un grup amb un valor representatiu; \emph{desviació típica} $\rightarrow$
        mesurar la desigualtat/dispersió; \emph{piràmide/histograma} $\rightarrow$
        veure la forma d'una població o distribució; \emph{lectura crítica}
        $\rightarrow$ no deixar-se enganyar per dades mal presentades; \emph{enquesta i
        buidatge} $\rightarrow$ recollir dades pròpies i convertir-les en informació.
  \item Resposta oberta. Exemple: «He après a \emph{descriure} un grup amb paràmetres,
        a \emph{visualitzar-lo} amb gràfics, a \emph{llegir} dades amb mirada crítica i
        a \emph{recollir-ne} de pròpies amb una enquesta i el full de càlcul.»
\end{sollist}
\end{solucions}
